\chapter{线性代数}

\section{行列式抓结构}
\recognize{问高阶行列式、含参数行列式或某次幂的系数。先找零块、相同行列、成对对称项，再决定展开位置。}
\core{交换两行变号；一行乘$k$使行列式乘$k$；倍加不变。分块三角阵满足$\left|\begin{smallmatrix}B&C\\0&D\end{smallmatrix}\right|=|B||D|$。行列式对每一行分别线性，但$|A+B|$通常不等于$|A|+|B|$。}
\process{\begin{enumerate}\item 用倍加制造零，或同时重排行、列暴露小块。\item 记录交换次数和提出的因子；沿零最多的行列展开。\item 求系数时只保留能产生目标次数的排列，不展开整个多项式。\item 用特殊参数值检验符号与次数。\end{enumerate}}
\fast{看到四阶阵中每行只有两个非零元，先试把奇偶行列分组。若变成$\left(\begin{smallmatrix}0&B\\C&0\end{smallmatrix}\right)$且每块为$m$阶，行列式为$(-1)^m|B||C|$。盲目按第一行展开会重复计算相似的余子式。}
\budget{易至中；选择、填空2--4分钟。目标系数题先用1分钟分析取项。}
\fullquestion{math2-2014-07}
\apply{题面零元呈交错分布，应先重排为两个二阶块。所得因子均为$ad-bc$，再核对交换的负号，结果为$-(ad-bc)^2$，选B。取$a=d=1,b=c=0$时原阵是一次交换的置换阵，行列式为$-1$，能立即排除正号。}
\related{2020-14（对称四阶阵）；2021-16（指定次数系数）；2023-16（行列式间的变换）。}

\section{余子式合并}
\recognize{题目直接给出$A_{ij}$的和、差，或把元素与其代数余子式联系起来。把这一串系数看成一次按行、列展开。}
\core{$\sum_j a_{ij}A_{ij}=|A|$；$\sum_j a_{kj}A_{ij}=0$（$k\ne i$）。固定第$i$行的余子式组合$\sum_j c_jA_{ij}$，等于把第$i$行换为$(c_1,\ldots,c_n)$后的行列式。$A^*=(A_{ji})$，余子式排成矩阵后还要转置才是伴随矩阵。}
\process{\begin{enumerate}\item 看余子式下标是否固定同一行或列。\item 将前面的系数补齐成替换行或列。\item 对新行列式做倍加、零块分解。\item 元素等式先改写成矩阵等式，再取行列式或乘原阵。\end{enumerate}}
\fast{只求两个余子式的差时，分别求两个三阶行列式往往多算一倍。固定下标的展开结构比数值更值得先看。}
\budget{中；2--4分钟。}
\fullquestion{math2-2013-14}
\apply{题面$a_{ij}+A_{ij}=0$对所有$i,j$成立，先写$A+(A^*)^T=0$，故$AA^T=-|A|E$。若$|A|=0$，则$AA^T=0$推出$A=0$，与非零条件矛盾；故可取行列式，得到$|A|^2=-|A|^3$，于是$|A|=-1$。关键链是“余子式整体关系$\to$伴随$\to$正半定矩阵”。}
\related{2019-14（余子式之差）；2012-14（交换行后与伴随相乘）。}

\section{伴随先分秩}
\recognize{出现$A^*$、某个余子式非零、齐次解的个数或$A(A-A^*)=0$。先判$A$可逆与否，再决定能否消去$A$。}
\core{$AA^*=A^*A=|A|E$。$n\ge2$时，$r(A^*)=n,1,0$分别对应$r(A)=n,n-1,\le n-2$。可逆时$A^*=|A|A^{-1}$；$|A^*|=|A|^{n-1}$。}
\process{\begin{enumerate}\item 用基础解系个数$n-r(A)$或非零余子式确定秩。\item 按$|A|\ne0$与$|A|=0$分支，化简伴随恒等式。\item 奇异分支用秩、核空间或矩阵多项式约束。\item 检查所有候选秩是否真能实现。\end{enumerate}}
\fast{遇到$A^*x=0$且$A$奇异，由$A^*A=0$知$A$的每一列都是解；若$r(A)=n-1$，挑出$n-1$个独立列便是基础解系。}
\budget{中；2--4分钟，含存在性验证约5分钟。}
\fullquestion{math2-2024-09}
\apply{由题设得$A^2=AA^*=|A|E$。若可逆，左乘$A^{-1}$将得到$A=A^*$，矛盾，故$|A|=0$，于是$A^2=0$。因此$\operatorname{Im}A\subseteq\ker A$，四阶时$r(A)\le2$；$A=0$又与$A\ne A^*$矛盾，故只可能1或2，两者均可由平方为零的矩阵实现，选D。}
\related{2019-07（由基础解系求伴随秩）；2020-07（伴随齐次解）；2023-08（分块伴随）。}

\section{变换辨清左右}
\recognize{问“若干初等行变换得到”、$PAQ$、置换矩阵或已知变换后的逆矩阵信息。第一动作是标记左乘改行、右乘改列。}
\core{行等价$B=PA$（$P$可逆）保持行空间和$\ker A$；列等价$B=AQ$保持列空间。矩阵等价$B=PAQ$只要求同型同秩。置换矩阵满足$P^{-1}=P^T$，$|P|=\pm1$，故$P^*=\pm P^{-1}$。}
\process{\begin{enumerate}\item 将题面变换写成乘积，保持操作先后次序。\item 行变换题比较行空间或齐次方程，不只比秩。\item 需恢复原阵时逆序撤销操作；求逆则按$(PAQ)^{-1}=Q^{-1}A^{-1}P^{-1}$。\item 选项题优先检查列间线性关系是否保持。\end{enumerate}}
\fast{只做行变换时，原矩阵各列之间的线性关系完全保留。这比把四个选项全部化阶梯形更快。}
\budget{易至中；2--3分钟。}
\fullquestion{math2-2025-09}
\apply{目标矩阵的列满足$c_2=c_1$、$c_4=c_1+2c_3$，且$c_1,c_3$独立。选项B保留这两条列关系且秩为2，故具有相同行空间，选B。这里只需比较两条关系；只检验秩会把“矩阵等价”误当成“行等价”。}
\related{2013-07（右乘可逆保持列空间）；2021-10（左右变换）；2022-16（逆序恢复）；2026-08（置换矩阵）。}

\section{矩阵方程分列}
\recognize{求$AX=B$、$XA=B$、右逆或“存在矩阵使等式成立”。把未知矩阵拆为未知列或未知行，立刻转成同一系数矩阵的多个方程组。}
\core{$AX=B$可解当且仅当$r(A,B)=r(A)$，即$B$每列均在$A$的列空间；解为$X=X_0+N$，其中$AN=0$。$XA=B$可转置为$A^TX^T=B^T$。$AC-CA=B$可逐元素设未知数，且必有$\operatorname{tr}B=0$。}
\process{\begin{enumerate}\item 写清各阵尺寸，确定应分列还是分行。\item 先用左零空间约束$B$，或化增广矩阵判可解。\item 统一求出齐次基础解系，再给每列加独立参数。\item 另要求$X$可逆时，最后补$|X|\ne0$。\end{enumerate}}
\fast{矩阵多列的可解条件往往共用一个左零向量$u$：$u^TA=0\Rightarrow u^TB=0$。这一步可绕过整套未知元。}
\budget{中；小题2--4分钟，大题8--12分钟。}
\fullquestion{math2-2026-09}
\apply{按完整题干的行列读取，$u=(0,2,-1,1)^T$满足$u^TA=0$。两列右端分别要求$2\cdot1-1+a=0$、$2\cdot1-1+b=0$，得到$a=b=-1$，选A。读到“存在$B$”时要先找相容条件，无需真的求$B$。}
\related{2014-22（$AB=E$的全部右逆）；2018-23（可逆解）；2013-22（交换子矩阵方程）。}

\section{含参方程分支}
\recognize{问唯一解、无解、无穷多解、参数取值。先算系数行列式寻找分支点，再在奇异点检查增广秩。}
\core{$r(A)=r(A,b)=n$唯一解；$r(A)=r(A,b)<n$无穷多解；$r(A)<r(A,b)$无解。行列式为零只说明不能唯一解，不能单独区分另两类。}
\process{\begin{enumerate}\item 求$|A|$或在不除含参式的前提下消元。\item 列出使主元消失的全部参数值。\item 对各奇异值单独代入，看最后零行右端是否为零。\item 写特解加齐次通解，参数数目核对$n-r(A)$。\end{enumerate}}
\fast{范德蒙德型列$(1,t,t^2)^T$出现时，参数相等即重复列；右端若也具有该结构，直接判断它是否落在相同列空间。}
\budget{中；选择3分钟，大题9--12分钟。}
\fullquestion{math2-2015-07}
\apply{三列对应节点$1,2,a$。$a\notin\{1,2\}$时可逆，只能唯一解；要无穷多解必须$a\in\{1,2\}$。此时右端$(1,d,d^2)^T$在前两列张成的平面上等价于$(d-1)(d-2)=0$，所以$d\in\{1,2\}$，选D。题面的平方结构让“范德蒙德+插值关系”比三参数硬消元更短。}
\related{2012-22（参数奇异与通解）；2022-09（范德蒙德可解分类）；2016-22（无解条件接正规方程）。}

\section{秩与零空间}
\recognize{给出列间关系、不同特征值、$r(AB)$与$r(BA)$，问解空间维数、公共非零解。把每个条件转换成秩上界、下界或核空间维数。}
\core{$\dim\ker A=n-r(A)$；$r(AB)\le\min\{r(A),r(B)\}$；$r(A)+r(B)-n\le r(AB)$。同列数矩阵的公共零空间为$\ker\left(\begin{smallmatrix}A\\B\end{smallmatrix}\right)$。$\dim(U\cap V)\ge\dim U+\dim V-n$。}
\process{\begin{enumerate}\item 从显式列关系给出上界，从独立向量或非零子式给出下界。\item 产品秩不等式先估计，不急着设矩阵全部元素。\item 公共解合并方程，或把两个核空间的维数相加。\item 存在性结论用维数严格大于零收尾。\end{enumerate}}
\fast{有$n$个不同特征值意味着可对角化，秩等于非零特征值的个数。已知一条列关系又只有一个零特征值时，秩便被夹定。}
\budget{中至难；3--5分钟。}
\fullquestion{math2-2025-10}
\apply{若$A$或$B$可逆，则$r(AB)=r(BA)$，故二者都奇异，$r(AB)\le2$，从而$r(BA)\le1$。又$r(ABA),r(BAB)\le r(BA)\le1$，两个核空间维数均至少2，在三维空间中必有非零交集，选D。这里“秩差1”应触发维数估计，逐项设矩阵无法有效控制参数。}
\related{2017-22（特征值与列关系夹秩）；2018-08（分块秩）；2016-14（矩阵等价即同秩）。}

\section{通解直接读取}
\recognize{题面已给列的独立性和一个线性关系，右端又是这些列的组合。直接读出特解和齐次方向。}
\core{$Ax=b$的全部解为$x=x_0+\sum_{i=1}^{n-r}k_i\eta_i$，其中$Ax_0=b$、$A\eta_i=0$。$\sum c_j\alpha_j=0$的系数向量$(c_j)$就是齐次解。}
\process{\begin{enumerate}\item 数独立列，确定秩和自由参数数。\item 从题面的列关系抄出齐次解。\item 从右端的列组合抄出一个特解。\item 用维数保证所得齐次向量已经覆盖全部解。\end{enumerate}}
\fast{有现成的列关系时，不必把符号列向量展开成坐标。特解并不唯一，选择最容易直接核验的那个即可。}
\budget{易；1--2分钟。}
\fullquestion{math2-2025-16}
\apply{前三列独立，所以$r(A)=3$、齐次解空间一维；$a_1+a_2-a_3-a_4=0$给出$\eta=(1,1,-1,-1)^T$。右端$a_1+4a_4$直接给$x_0=(1,0,0,4)^T$。因此$x=x_0+k\eta$。从题干读系数即可完成，不需要求任何行列式。}
\related{2017-22（抽象列向量的非齐次通解）；2014-22（每列分别添加齐次解）。}

\section{张成空间辨析}
\recognize{问线性表出、向量组等价、公共可表向量、任意两向量独立。先把文字译成子空间包含或交集。}
\core{组$A$可由组$B$表出等价于$\operatorname{Im}A\subseteq\operatorname{Im}B$，也等价于存在$C$使$A=BC$。两组等价要求张成空间相同，不能只看秩相同。$U\subseteq V\Rightarrow V^\perp\subseteq U^\perp$，而$(\operatorname{Im}A)^\perp=\ker A^T$。}
\process{\begin{enumerate}\item 包含问题写$A=BC$；公共表示写$As=Bt$。\item 同时消元$[A,-B]$，把系数解代回得到公共向量。\item 参数题在秩改变处另验空间是否相同。\item 三向量相关且两两独立时，相关关系中三个系数都非零。\end{enumerate}}
\fast{“列向量可表出”对应转置方程的解包含；看到$A^Tx=0$就想到与所有列正交。直接把结论套到$Ax=0$容易把对象颠倒。}
\budget{中；2--4分钟，含参多分支约5分钟。}
\fullquestion{math2-2021-09}
\apply{题面给出$\operatorname{Im}A\subseteq\operatorname{span}\{\beta_1,\beta_2\}\subseteq\operatorname{Im}B$。取正交补得$\ker B^T\subseteq\ker A^T$，所以选D。识别“列空间包含”后只需反向取正交补。}
\related{2022-10（参数与等价）；2023-10（两个张成空间之交）；2024-16（两两独立）；2019-22（等价与线性表示）。}

\section{对称乘积降维}
\recognize{出现$A^TA$、$AA^T$、平方和或正规方程。优先利用范数和秩，不急着展开乘积。}
\core{实矩阵满足$x^TA^TAx=\|Ax\|^2$，故$\ker(A^TA)=\ker A$、$r(A^TA)=r(A)$；$AA^T$也同秩，且均半正定。$A^TAx=A^Tb$总相容；其齐次部分仍为$\ker A$。}
\process{\begin{enumerate}\item 用规范形的非零项数确定乘积的秩。\item 改求原矩阵的秩条件，小阶时令相关行列成比例。\item 正规方程先求一个特解，齐次部分直接沿用$Ax=0$。\item 平方和为零时逐个线性式置零。\end{enumerate}}
\fast{$AA^T$尺寸可能更大，但其秩并未增加。目标只问秩或零空间时，乘开会把线性条件变成二次条件。}
\budget{易至中；2--3分钟。}
\fullquestion{math2-2026-16}
\apply{规范形只有一个正平方项，故$r(AA^T)=r(A)=1$。按题面$A=\left(\begin{smallmatrix}1&b&-1\\a+2&3&-3a\end{smallmatrix}\right)$，第一、第三列子式为$2-2a$，得$a=1$；第一、第二列子式为$3-3b$，得$b=1$。所以$a+b=2$。识别“Gram矩阵的秩”即可避开$AA^T$的二次计算。}
\related{2016-22（无解方程与正规方程）；2018-22（平方和）；2013-23（向量外积二次型）。}

\section{先算作用再求谱}
\recognize{题面给$A\alpha_i$的表达、$P^{-1}AP$、线性变换作用式，或特征向量含参数。先从“作用到基向量后是什么”提取坐标矩阵。}
\core{$AP=PC$时，$C$第$j$列是$A\alpha_j$在基$P=(\alpha_1,\ldots,\alpha_n)$下的坐标；$P$可逆则$C=P^{-1}AP$。特征向量满足$A\alpha=\lambda\alpha$，$\alpha\ne0$。$f(A)$的特征值为$f(\lambda_i)$。}
\process{\begin{enumerate}\item 依次读取每个$A\alpha_j$，把系数竖着放为第$j$列。\item 在坐标矩阵中算谱、行列式或幂。\item 已知一个特征向量时先比较最简单分量求$\lambda$，再解参数。\item 检查特征值的和等于迹、积等于行列式。\end{enumerate}}
\fast{换基后的小矩阵往往已分块三角。不要反求原矩阵$A$；要求的谱和行列式本来就不随相似变化。}
\budget{易至中；2--4分钟。}
\fullquestion{math2-2018-14}
\apply{三个给定向量独立，直接构造坐标矩阵$C=\left(\begin{smallmatrix}2&0&0\\1&1&-1\\1&2&1\end{smallmatrix}\right)$。右下块特征方程为$(1-\lambda)^2+2=0$，故唯一实特征值为2。题干中三个“$A\alpha=$”就是按列构造坐标矩阵的信号。}
\related{2017-14（单个特征向量求参数）；2015-14（矩阵多项式的行列式）；2023-22（作用式恢复矩阵）；2020-23（循环基）。}

\section{重根查维数}
\recognize{问相似于对角矩阵、仅有两个不同特征值、含参相似。先求特征值及重数，再检查重根拥有多少独立特征向量。}
\core{$A$可对角化当且仅当各特征空间维数之和为$n$。重数为$m$的特征值$\lambda$必须满足$n-r(A-\lambda E)=m$。同谱是相似的必要条件；可对角化矩阵同谱（计重数）才可据此判相似。}
\process{\begin{enumerate}\item 利用三角或分块结构分解特征多项式。\item 用“仅两个不同特征值”列出重合分支。\item 在每个分支计算$r(A-\lambda E)$，筛选参数。\item 将特征向量按目标对角元的顺序放入$P$。\end{enumerate}}
\fast{单根自动贡献一个维数，工作集中在重根处。看到重根不要直接宣布不可对角化，也不必计算每个根的全部向量后才检查。}
\budget{中；大题10--14分钟。}
\fullquestion{math2-2021-22}
\apply{分块三角结构直接给特征值$1,3,b$，故$b=1$或3。$b=1$时要求$r(A-E)=1$，得到$a=1$；$b=3$时要求$r(A-3E)=1$，得到$a=-1$。接下来只在对应分支求两维重根特征空间与一维单根空间。题面“仅两个”负责找分支，“相似于对角阵”负责筛参数。}
\related{2017-08（同谱与重根）；2019-23（相似不变量后求$P$）；2013-08（对称矩阵与相似条件）。}

\section{共同特征空间}
\recognize{出现$AB=BA$、调整对角化矩阵的列、同一特征值的多个特征向量。抓住特征空间在另一个矩阵作用下保持不变。}
\core{$A\alpha=\lambda\alpha$且$AB=BA$时，$A(B\alpha)=\lambda B\alpha$。若$\lambda$为单特征值，其特征空间一维，故$B\alpha=\mu\alpha$。同一特征值的特征向量可做可逆线性组合；不同特征值的向量一般不能混加成特征向量。}
\process{\begin{enumerate}\item 先判特征空间是一维还是高维。\item 交换条件将每个特征空间限制为不变子空间。\item 单根情形沿用$A$的特征向量组成$B$的特征基。\item 判断必要性时取$A=E$或$B=E$做具体检验。\end{enumerate}}
\fast{“$A$有互异特征值”是强结构，能固定各条特征方向；“$A$可对角化”在重根时没有这么强。}
\budget{中；2--4分钟。}
\fullquestion{math2-2024-10}
\apply{二阶$A$有两个不同特征值，提供两条一维特征空间；交换条件迫使$B$保持这两条方向，所以$B$可对角化。反过来取$B=E,A=E$，结论成立而前提不成立，因此是充分不必要条件，选B。}
\related{2012-08（同特征空间换基）；2020-08（特征向量排序）；2016-07（相似与转置运算）。}

\section{多项式压缩矩阵}
\recognize{给$A^3=0$、$(A-B)^2=0$、$A^2+A=2E$或复杂矩阵方程。先把等式整理成一个矩阵的低次多项式。}
\core{$N^m=0\Rightarrow(E-N)^{-1}=E+N+\cdots+N^{m-1}$。$p(A)=0$则每个特征值满足$p(\lambda)=0$。若$N^2=0$，则$\operatorname{Im}N\subseteq\ker N$，$r(N)\le n/2$。实对称矩阵可正交对角化，因此对称且幂零必为零阵。}
\process{\begin{enumerate}\item 把两阵等式尝试写成$(A-B)^2$，同时保留乘法次序。\item 用多项式限制谱，用秩限制特征空间维数。\item 要逆矩阵时用有限几何级数；要解矩阵方程时先因式分解。\item 检查$A\ne B$等非零条件，排除退化分支。\end{enumerate}}
\fast{矩阵不能任意交换，但$(A-B)^2=A^2-AB-BA+B^2$始终成立。整理成一个新矩阵后，秩与谱往往一起变得简单。}
\budget{中；小题2--4分钟，大题8--12分钟。}
\fullquestion{math2-2026-10}
\apply{置$N=A-B\ne0$，题设即$N^2=0$。三阶时$r(N)\le1$且非零，故$r(N)=1$、$\dim\ker N=2$；它只有零特征值，却有两个独立特征向量，D错误。若$A,B$均对称，则$N$对称且幂零，只能为零，也与题设矛盾。}
\related{2015-22（幂零与矩阵方程因式分解）；2019-08（多项式限定二次型惯性）。}

\section{矩阵幂先降阶}
\recognize{问$A^{99}$、$A^{10}$，且列向量高度相关、给出极大无关组，或存在$A=GH$。先寻找低维矩阵或低次递推。}
\core{$A=P\Lambda P^{-1}\Rightarrow A^m=P\Lambda^mP^{-1}$。$A=GH$时$A^m=G(HG)^{m-1}H$（$m\ge1$），其中$HG$可能远小于$A$。若$HG=E+N,N^2=0$，则$(HG)^k=E+kN$。}
\process{\begin{enumerate}\item 用极大无关组组成$G$，把每列的表示系数组成$H$。\item 计算较小的$HG$，识别对角、幂零或递推结构。\item 算完低维幂后乘回，保留因子形式或按要求展开。\item 用$m=1,2$检验指数位置。\end{enumerate}}
\fast{给“极大无关组$\to A=GH\to$求幂”的连续问法，第一问就在给第二问铺路。直接对四阶阵求特征向量会失去题面提供的降维优势。}
\budget{中；10--14分钟，低秩分解通常先用3分钟完成。}
\fullquestion{math2-2026-22}
\apply{$\alpha_3=-\alpha_1+\alpha_2$，$\alpha_4=\alpha_1-\alpha_2$，而$\alpha_1,\alpha_2$独立，故为极大无关组。依次读系数得
\[H=\begin{pmatrix}1&0&-1&1\\0&1&1&-1\end{pmatrix},\qquad HG=\begin{pmatrix}1&-1\\0&1\end{pmatrix}=E+N.\]
因$N^2=0$，$A^{10}=G\left(\begin{smallmatrix}1&-9\\0&1\end{smallmatrix}\right)H$。这里指数是9，因为两端$G,H$已经贡献一次$A$。}
\related{2016-23（高次幂与列向量表示）；2014-23（秩一矩阵相似）；2020-23（循环基求坐标矩阵）。}

\section{二次型三种变换}
\recognize{题目交替出现“合同”“正交变换”“标准形”“规范形”。先写出允许的变换，再决定哪些数值必须保留。}
\core{$x=Cy$给$B=C^TAC$，称合同，保持秩与正、负惯性指数；规范形系数为$1,-1,0$。正交变换满足$Q^TQ=E$，此时$Q^TAQ=Q^{-1}AQ$，对角元是原矩阵特征值。一般标准形的非零系数不要求为$\pm1$。}
\process{\begin{enumerate}\item 由题面二次项写对称阵，交叉项系数除以2。\item 合同条件只比较正负个数与零个数。\item 正交条件还要匹配全部特征值，可用迹与行列式。\item 求正交阵时分别求特征空间，再单位化。\end{enumerate}}
\fast{合同题只问参数范围时，先数符号；正交题若目标对角阵含参数，先用迹求参数。无须过早解全部特征向量。}
\budget{中；小题2--4分钟，大题12--16分钟。}
\fullquestion{math2-2025-22}
\apply{目标$B=\operatorname{diag}(k,6,0)$有零特征值，故$|A|=-3(a-4)=0$，得$a=4$。此时$A$的特征值为$3,6,0$，惯性为$(2,0)$，故$k>0$。第二问添加“正交”，由$k+6=\operatorname{tr}A=9$立刻得$k=3$。分别取单位特征向量$(1,1,1)^T/\sqrt3$、$(-1,0,1)^T/\sqrt2$、$(1,-2,1)^T/\sqrt6$作为$Q$的三列。第一问保符号，第二问保特征值。}
\related{2020-22（一般可逆变换）；2015-08（正交列重排）；2022-08（等价、相似、合同辨析）。}

\section{惯性看配方}
\recognize{二次型本来写成平方和、平方差，或只问正负惯性指数和规范形。先看线性式是否独立，尽量不求三次特征方程。}
\core{实对称阵的惯性等于正、负、零特征值的个数。若$f=\|Cx\|^2$，则只有正与零项，正惯性指数为$r(C)$。配方中每次使用的变量替换应可逆；不能把线性相关的平方项直接当独立新变量。}
\process{\begin{enumerate}\item 将平方中的线性式排成系数阵$C$，先查秩。\item 独立则直接作为新变量；相关则先消去一个线性式。\item 配成平方差，数正负项，核对总秩。\item 含参时在配方系数为零处单独讨论。\end{enumerate}}
\fast{没有平方项而只有$2xy$时用$u=x+y,v=x-y$，有$2xy=(u^2-v^2)/2$。只求惯性不必把变换矩阵单位化。}
\budget{易至中；2--4分钟。}
\fullquestion{math2-2021-08}
\apply{令$u=x_1+x_2,v=x_2+x_3$，则$x_3-x_1=v-u$，故$f=u^2+v^2-(v-u)^2=2uv$。取第三个独立线性式补成可逆变换，再令$p=u+v,q=u-v$，得$f=(p^2-q^2)/2$。正、负惯性指数均为1，选C。三个平方并不独立，是读题时最关键的识别点。}
\related{2018-22（平方和随参数变秩）；2023-09（平方差与相关线性式）；2016-08（特征值符号区间）；2014-14（负惯性指数）。}

\section{正定与商的最值}
\recognize{问二次型正定、特征值正负个数，或$\min f(x)/(x^Tx)$。先用正交对角化把问题变成带权平均。}
\core{实对称$A$正定等价于全部特征值正，也等价于全部顺序主子式正。$\lambda_{\min}\le x^TAx/(x^Tx)\le\lambda_{\max}$（$x\ne0$），等号恰在相应特征空间取得。半正定要检查全部主子式非负，不能只照搬顺序主子式条件。}
\process{\begin{enumerate}\item 按结构拆成小块，或写成$aE+b\mathbf1\mathbf1^T$。\item 只问正负时用行列式、迹、主子式判别，避免求根。\item 商的最值求最小、最大特征值及对应特征向量。\item 重根的特征空间内做正交化，最后单位化。\end{enumerate}}
\fast{二阶实对称块若行列式负，必一正一负；再看独立的一阶块即可数三阶阵惯性。商中分母为$x^Tx$时，不必用拉格朗日乘数重新计算。}
\budget{中；小题2--3分钟，大题10--14分钟。}
\fullquestion{math2-2022-22}
\apply{矩阵为$\left(\begin{smallmatrix}3&0&1\\0&4&0\\1&0&3\end{smallmatrix}\right)$。$x_1,x_3$对称块给特征值$2,4$，另一个方向也给4；单位特征向量可取$(1,0,-1)^T/\sqrt2,(0,1,0)^T,(1,0,1)^T/\sqrt2$。变换后$f=2y_1^2+4y_2^2+4y_3^2$，所以商的最小值2，在$y_2=y_3=0$时取得。}
\related{2025-08（二阶块符号）；2019-08（多项式、行列式与惯性）；2013-23（正交单位向量与谱）。}

\section{线代读题路线}
\begin{tabularx}{\textwidth}{p{31mm}Y Y}\toprule
题面结构&第一动作&接着串联的知识\\\midrule
余子式、伴随&写$AA^*=|A|E$并分秩&可逆、零空间、幂零\\
列关系、极大无关组&读系数，定秩&特解、基础解系、低秩分解\\
存在$X$使$AX=B$&查右端列是否相容&左零空间、增广秩\\
含参数且可对角化&找重根，算特征空间维数&秩、相似、特征向量排列\\
平方和、$A^TA$&改看原矩阵的秩&半正定、零空间、规范形\\
合同后再加正交&先惯性，再特征值&行列式、迹、正交化\\
高次幂且列相关&构造$A=GH$&低维幂、幂零二项式\\\bottomrule
\end{tabularx}

\begin{card}{teal}{最后两分钟检查}
\begin{itemize}
\item $P$的列顺序是否对应对角元顺序；正交阵是否已单位化，重根空间内是否正交。
\item 通解参数个数是否等于未知数个数减秩；特解是否真的代回右端。
\item 含参消元有没有除掉可能为零的主元；相同秩有没有被误当相同行空间。
\item 二次型交叉系数是否除以2；合同有没有误用特征值保持；$GH$的高次幂指数是否少1。
\end{itemize}
\end{card}
